\section{Aprendizaje autónomo: la gradación que ya está ocurriendo}

\subsection{Los múltiples sentidos del aprendizaje autónomo}

El «aprendizaje autónomo» se ha vuelto un tema candente en Silicon Valley, pero cada persona lo define de manera distinta. Yao Shunyu (姚顺雨) lo divide en varios tipos:

\begin{itemize}
\item \textbf{Personalización del chat}: ChatGPT aprovecha los datos del usuario para aprender continuamente el estilo de conversación
\item \textbf{Adaptación al entorno de código}: se familiariza cada vez más con el entorno y la documentación particulares de cada empresa
\item \textbf{Exploración científica}: aprender química orgánica desde cero como un estudiante de doctorado hasta convertirse en experto del área
\end{itemize}

\begin{importantbox}{Observación clave: el cuello de botella está en los datos y las tareas, no en la metodología}
El verdadero cuello de botella del aprendizaje autónomo no es la metodología (how), sino los datos y las tareas (what and where): en qué escenario, bajo qué premisa y con qué función de recompensa se realiza.
\end{importantbox}

\subsection{El aprendizaje autónomo ya está ocurriendo}

\begin{knowledgebox}{La automejora de Claude Code}
Claude Code ya ha escrito el 95\% del código de su propio proyecto; en cierto sentido, se está ayudando a sí mismo a mejorar. ¿No es esto una forma de aprendizaje autónomo?

De manera similar, el AutoComplete Model de Cursor se entrena cada pocas horas con los datos más recientes de los usuarios, y el Composer Model también usa datos de entornos reales. Todas estas son señales tempranas de aprendizaje autónomo.
\end{knowledgebox}

Cuando Yao Shunyu (姚顺雨) trabajaba en SWE-Agent en 2022--2023, ya planteaba: si algún día SWE-Agent pudiera mejorar por sí mismo el repositorio de código de SWE-Agent, eso sería una forma de aprendizaje autónomo. Claude Code ya está haciendo esto a gran escala.

\begin{warningbox}{Gradación, no salto abrupto}
Yao Shunyu (姚顺雨) considera que el aprendizaje autónomo se parece más a una gradación que a un salto abrupto: los ejemplos actuales todavía están limitados a escenarios específicos, pero la tendencia es clara. El mayor cuello de botella es la \textbf{imaginación}: si en 2026--2027 surge un nuevo paradigma, ¿cómo debería verse esa entrada de blog? ¿Qué tipo de tareas y resultados debería mostrar?
\end{warningbox}

